\documentclass[11pt,a4paper]{article}

\usepackage[utf8]{inputenc}
\usepackage[T1]{fontenc}
\usepackage{graphicx}
\usepackage{booktabs}
\usepackage{hyperref}
\usepackage{listings}
\usepackage{xcolor}
\usepackage{geometry}
\usepackage{amsmath}
\usepackage{float}

\geometry{margin=1in}

\definecolor{codegreen}{rgb}{0,0.6,0}
\definecolor{codegray}{rgb}{0.5,0.5,0.5}
\definecolor{codepurple}{rgb}{0.58,0,0.82}
\definecolor{backcolour}{rgb}{0.95,0.95,0.92}

\lstdefinestyle{mystyle}{
    backgroundcolor=\color{backcolour},
    commentstyle=\color{codegreen},
    keywordstyle=\color{magenta},
    numberstyle=\tiny\color{codegray},
    stringstyle=\color{codepurple},
    basicstyle=\ttfamily\footnotesize,
    breakatwhitespace=false,
    breaklines=true,
    captionpos=b,
    keepspaces=true,
    numbers=left,
    numbersep=5pt,
    showspaces=false,
    showstringspaces=false,
    showtabs=false,
    tabsize=2
}
\lstset{style=mystyle}

\title{Evaluating Detection Rule Efficacy: A Reproducible Benchmark for Security Operations}
\author{SOC Detection Engineering Benchmark}
\date{\today}

\begin{document}

\maketitle

\begin{abstract}
Security Operations Centers rely on detection rules to identify malicious activity, yet systematic evaluation of rule efficacy remains challenging. We present a reproducible benchmark framework that measures detection rule performance against simulated attack behaviors mapped to MITRE ATT\&CK. Using a containerized lab environment with Atomic Red Team tests and Sigma detection rules, we evaluate true positive rate, false positive rate, mean time to detect, and alert volume across 11 ATT\&CK techniques. Our results reveal significant gaps in detection coverage for specific techniques, with an overall detection rate of \textbf{[X]\%} and a false positive rate of \textbf{[Y]} alerts per hour. We provide error analysis explaining detection failures and release all artifacts for reproducibility.
\end{abstract}

\section{Introduction}

Detection engineering is a critical discipline within security operations, yet practitioners lack standardized methods for evaluating detection rule effectiveness. Questions such as ``What percentage of attacks do our rules detect?'' and ``How many false positives do our rules generate?'' often go unanswered due to the absence of systematic benchmarking approaches.

This paper presents a reproducible framework for evaluating detection rules against simulated attack behaviors. Our contributions include:

\begin{itemize}
    \item A containerized lab environment for reproducible experiments
    \item A benchmark harness that executes ATT\&CK-mapped attacks and measures detection metrics
    \item Empirical evaluation of Sigma rules against 11 techniques across multiple tactics
    \item Error analysis explaining detection gaps and their root causes
    \item Fully open-source artifacts for reproducibility
\end{itemize}

\section{Background and Related Work}

\subsection{MITRE ATT\&CK Framework}

The ATT\&CK framework provides a taxonomy of adversary tactics and techniques based on real-world observations. We use ATT\&CK technique identifiers (e.g., T1059.004) to map attacks to detection rules, enabling systematic coverage analysis.

\subsection{Sigma Detection Rules}

Sigma is a generic signature format for SIEM systems that enables portable detection rules. Rules are written in YAML and compiled to target-specific query languages (e.g., Elasticsearch Query DSL, Splunk SPL).

\subsection{Atomic Red Team}

Atomic Red Team provides small, portable tests mapped to ATT\&CK techniques. Each ``atomic'' test simulates a specific adversary behavior with minimal dependencies, enabling automated execution.

\subsection{Related Work}

Previous work on detection evaluation includes:
\begin{itemize}
    \item MITRE ATT\&CK Evaluations for endpoint detection products
    \item Purple team exercises combining red and blue team activities
    \item Detection-as-code practices in DevSecOps
\end{itemize}

Our work differs by providing a fully automated, reproducible benchmark with quantitative metrics.

\section{Threat Model}

We scope our evaluation to a Linux-based enterprise environment with the following assumptions:

\begin{itemize}
    \item \textbf{Attacker capabilities}: Post-compromise attacker with shell access
    \item \textbf{Attack scope}: Execution, Persistence, Privilege Escalation, Defense Evasion, Discovery, and Command \& Control tactics
    \item \textbf{Defender capabilities}: Host-based logging (auditd) with centralized SIEM (Elasticsearch)
    \item \textbf{Out of scope}: Initial access, lateral movement, data exfiltration
\end{itemize}

\section{Methodology}

\subsection{Lab Architecture}

Our benchmark uses a Docker Compose stack with four components:

\begin{enumerate}
    \item \textbf{Target Container}: Ubuntu 22.04 with auditd logging and Atomic Red Team
    \item \textbf{Log Shipper}: Filebeat shipping logs to Elasticsearch
    \item \textbf{Detection Engine}: Elasticsearch with compiled Sigma rules
    \item \textbf{Benchmark Harness}: Python orchestration for test execution and metrics collection
\end{enumerate}

\begin{figure}[H]
    \centering
    % \includegraphics[width=0.8\textwidth]{figures/architecture.png}
    \fbox{\parbox{0.8\textwidth}{\centering Architecture Diagram\\(Target $\rightarrow$ Filebeat $\rightarrow$ Elasticsearch $\rightarrow$ Detection)}}
    \caption{Lab architecture showing log flow from target to detection engine}
    \label{fig:architecture}
\end{figure}

\subsection{Attack Selection}

We selected 11 ATT\&CK techniques based on:
\begin{itemize}
    \item Prevalence in real-world incidents (MITRE data)
    \item Availability of Atomic Red Team tests for Linux
    \item Detectability via host-based logging
\end{itemize}

Table~\ref{tab:techniques} lists the techniques under evaluation.

\begin{table}[H]
    \centering
    \caption{ATT\&CK Techniques Under Evaluation}
    \label{tab:techniques}
    \begin{tabular}{lll}
        \toprule
        \textbf{ID} & \textbf{Technique} & \textbf{Tactic} \\
        \midrule
        T1059.004 & Unix Shell & Execution \\
        T1053.003 & Cron & Persistence \\
        T1136.001 & Local Account & Persistence \\
        T1548.001 & Setuid/Setgid & Privilege Escalation \\
        T1070.004 & File Deletion & Defense Evasion \\
        T1070.003 & Clear Command History & Defense Evasion \\
        T1027 & Obfuscated Files & Defense Evasion \\
        T1018 & Remote System Discovery & Discovery \\
        T1082 & System Information Discovery & Discovery \\
        T1105 & Ingress Tool Transfer & Command \& Control \\
        T1571 & Non-Standard Port & Command \& Control \\
        \bottomrule
    \end{tabular}
\end{table}

\subsection{Metrics}

We measure the following metrics:

\begin{itemize}
    \item \textbf{True Positive Rate (TPR)}: $\frac{\text{Attacks Detected}}{\text{Total Attacks}}$
    \item \textbf{False Negative Rate (FNR)}: $1 - \text{TPR}$
    \item \textbf{False Positive Rate (FPR)}: Alerts per hour during benign baseline
    \item \textbf{Mean Time to Detect (MTTD)}: Average latency from attack to alert
    \item \textbf{Alert Volume}: Average alerts per attack execution
\end{itemize}

\subsection{Experimental Procedure}

\begin{enumerate}
    \item \textbf{Baseline Phase}: Execute benign workloads for 10 iterations to establish FP baseline
    \item \textbf{Attack Phase}: Execute each atomic test with 2-second delay between attacks
    \item \textbf{Collection Phase}: Query Elasticsearch for alerts within attack time windows
    \item \textbf{Analysis Phase}: Calculate metrics and generate reports
\end{enumerate}

\section{Results}

% Replace [X] placeholders with actual results after running benchmark

\subsection{Overall Detection Performance}

\begin{table}[H]
    \centering
    \caption{Aggregate Benchmark Results}
    \begin{tabular}{lr}
        \toprule
        \textbf{Metric} & \textbf{Value} \\
        \midrule
        Total Attacks Executed & [N] \\
        Attacks Detected & [M] \\
        Overall Detection Rate & [X]\% \\
        Mean Time to Detect & [Y]s \\
        False Positive Rate & [Z]/hour \\
        Alert Volume per Attack & [V] \\
        \bottomrule
    \end{tabular}
\end{table}

\subsection{Detection by Technique}

Table~\ref{tab:results} shows per-technique detection rates.

\begin{table}[H]
    \centering
    \caption{Detection Rate by ATT\&CK Technique}
    \label{tab:results}
    \begin{tabular}{llrr}
        \toprule
        \textbf{Technique ID} & \textbf{Tactic} & \textbf{Detection Rate} & \textbf{MTTD (s)} \\
        \midrule
        T1059.004 & Execution & [X]\% & [Y] \\
        T1053.003 & Persistence & [X]\% & [Y] \\
        T1136.001 & Persistence & [X]\% & [Y] \\
        T1548.001 & Privilege Escalation & [X]\% & [Y] \\
        T1070.004 & Defense Evasion & [X]\% & [Y] \\
        T1070.003 & Defense Evasion & [X]\% & [Y] \\
        T1027 & Defense Evasion & [X]\% & [Y] \\
        T1018 & Discovery & [X]\% & [Y] \\
        T1082 & Discovery & [X]\% & [Y] \\
        T1105 & Command \& Control & [X]\% & [Y] \\
        T1571 & Command \& Control & [X]\% & [Y] \\
        \bottomrule
    \end{tabular}
\end{table}

\subsection{Detection Gaps}

The following techniques had 0\% detection rate:

\begin{itemize}
    \item [List techniques with detection failures]
\end{itemize}

\section{Error Analysis}

We analyzed detection failures to understand their root causes.

\subsection{Missing Log Coverage}

Some techniques failed detection due to insufficient audit rule coverage:

\begin{itemize}
    \item \textbf{[Technique]}: Auditd rules did not capture [specific syscalls/events]
\end{itemize}

\subsection{Rule Logic Gaps}

Detection rules contained logic gaps that prevented matching:

\begin{itemize}
    \item \textbf{[Technique]}: Sigma rule expected [condition] but attack produced [actual behavior]
\end{itemize}

\subsection{Log Parsing Issues}

Some detections failed due to log parsing inconsistencies:

\begin{itemize}
    \item \textbf{[Technique]}: Field name mismatch between auditd output and Sigma rule expectation
\end{itemize}

\section{Discussion}

\subsection{Implications for Detection Engineering}

Our results suggest that:
\begin{itemize}
    \item Detection rules require regular testing against real attack simulations
    \item Audit logging configuration significantly impacts detection capability
    \item False positive rates vary widely across technique categories
\end{itemize}

\subsection{Limitations}

\begin{itemize}
    \item Linux-only environment (no Windows coverage)
    \item Container-based target may differ from production environments
    \item Limited to auditd-based logging
    \item Single SIEM backend (Elasticsearch)
\end{itemize}

\subsection{Future Work}

\begin{itemize}
    \item Extend to Windows targets with Sysmon logging
    \item Add network-based detection (Suricata/Zeek)
    \item Integrate with commercial SIEM products
    \item Expand ATT\&CK technique coverage
\end{itemize}

\section{Reproducibility}

All artifacts are available at:

\begin{lstlisting}[language=bash]
# Clone the repository
git clone https://github.com/[org]/soc-detection-benchmark.git
cd soc-detection-benchmark

# Install dependencies
make install

# Start the lab
make up

# Run full benchmark
make benchmark

# View results
cat results/report_latest.md
\end{lstlisting}

\section{Conclusion}

We presented a reproducible benchmark for evaluating detection rule efficacy against simulated attacks. Our framework enables systematic measurement of true positive rate, false positive rate, and mean time to detect across ATT\&CK techniques. Error analysis revealed specific gaps in detection coverage that can inform rule improvement efforts. We release all artifacts to support reproducibility and community contribution.

\bibliographystyle{plain}
% \bibliography{references}

\appendix

\section{Sigma Rule Examples}

\begin{lstlisting}[language=yaml,caption={Example Sigma Rule for Unix Shell Execution}]
title: Linux Shell Execution
id: 9a5b8c3d-1234-5678-9abc-def012345678
status: experimental
logsource:
    product: linux
    service: auditd
detection:
    selection:
        type: EXECVE
        a0|endswith:
            - '/bash'
            - '/sh'
    condition: selection
level: low
\end{lstlisting}

\section{Attack Manifest Schema}

\begin{lstlisting}[language=yaml,caption={Attack Manifest Structure}]
techniques:
  - id: T1059.004
    name: "Unix Shell"
    tactic: execution
    atomic_tests:
      - name: "Execute shell commands"
        command: |
          /bin/bash -c 'whoami'
        cleanup: null
        expected_logs:
          - type: auditd
            key: exec
\end{lstlisting}

\end{document}
